This paper proposes the use of dynamic factor models to construct country-specific financial cycle indicators, capturing the common dynamics of variables related to credit, asset prices, indebtedness, and financial conditions. The proposed indicator explains a substantial share of the variance in these variables, effectively representing key aspects of financial cycle fluctuations.

Our findings show that the indicator provides a coherent narrative of the financial cycle's evolution across Europe, aligning closely with historical episodes of financial expansion and contraction. Moreover, it demonstrates strong early-warning capabilities, particularly in signalling heightened systemic risk in the lead-up to the GFC. In a simulated out-of-sample exercise, the estimated probability of a systemic banking crisis within one year more than doubles in the period immediately before the GFC, rising from approximately 10\% in early 2006 to around 20\% by 2007.

We also introduce a methodology to decompose the indicator into contributions from individual variables and define data-driven thresholds for identifying periods of neutral systemic risk. These tools offer valuable insights for monitoring cyclical systemic risk and informing macroprudential policy, particularly for the calibration of the countercyclical capital buffer.

% TODO (pre-submission): Expand the conclusion to 5-7 paragraphs covering:
%   1. Limitations (pseudo real-time assumptions: no vintages, no publication lags)
%   2. Directions for future research (non-European countries, dynamic loadings, nowcasting)
%   3. Broader policy takeaways (simplicity and transparency as features)
%   4. Connection to positive neutral CCyB literature
